\subsection{Actual polynomial inflation and the fibre annulus Gauss skeleton}
\begin{theorem}[Whole polynomial Gauss inflation]
\label{analytic:actual-polynomial-gauss-inflation}
\lean{EnvelopeProperty.Berkovich.RelativeSpectrum.centeredGauss_eq_gauss_of_val_le,
EnvelopeProperty.Berkovich.RelativeSpectrum.polynomialGaussInflation,
EnvelopeProperty.Berkovich.RelativeSpectrum.polynomialGaussInflation_val_apply,
EnvelopeProperty.Berkovich.RelativeSpectrum.continuous_polynomialGaussInflation_prod,
EnvelopeProperty.Berkovich.RelativeSpectrum.polynomialGaussInflation_zero,
EnvelopeProperty.Berkovich.RelativeSpectrum.proj_polynomialGaussInflation,
EnvelopeProperty.Berkovich.RelativeSpectrum.polynomialGaussInflation_val_X,
EnvelopeProperty.Berkovich.RelativeSpectrum.val_le_polynomialGaussInflation,
EnvelopeProperty.Berkovich.RelativeSpectrum.polynomialGaussInflation_eq_gauss,
EnvelopeProperty.Berkovich.RelativeSpectrum.polynomialGaussInflation_gauss}
\leanok
\uses{spectrum:gauss-maximality,spectrum:gauss-joint-continuity,spectrum:centered-gauss-whole-ray,spectrum:actual-line-gauss-inflation}
Let $R,A$ be arbitrary commutative rings, let $R\to A$ be the original algebra
map, and fix a real nonnegative valuation $v$ on $R$. For an ORIGINAL point
$x\in\operatorname{Sp}_v(A[T])$ and $s\geq0$, define $I_s(x)$ by taking the
Gauss valuation of radius $s$ over the WHOLE point $x$, translating the new
variable $U$ to $U+T$, and restricting along the original map
$A[T]\to A[T][U]$, $p\mapsto p(U)$ with every original coefficient retained.
Thus $|p(I_s(x))|$ is the genuine Gauss value of $p(U+T)$ with coefficients
evaluated by $x$. The construction is jointly continuous in $(x,s)$,
retains every original $A$ coefficient, and satisfies
\[
 I_0(x)=x,\qquad |T(I_s(x))|=\max(|T(x)|,s),\qquad
 |p(x)|\leq |p(I_s(x))| \quad(p\in A[T]).
\]
If $|T(x)|\leq s$, then $I_s(x)$ is the ENTIRE Gauss point of radius $s$
over the restriction of $x$ to $A$. If $x=\eta_t(y)$ is itself an original
Gauss point and $s\leq t$, then $I_s(x)=x$ as whole points.
These assertions require no field, completeness, reducedness, bounded
coefficient presentation, or nonzero radius hypothesis.
\end{theorem}
\begin{proof}
All maps in the construction are the actual polynomial algebra maps. Joint
continuity follows from the jointly continuous original Gauss construction
and the continuous restriction along a fixed ring homomorphism. Translation
fixes every constant, so the original coefficient restriction is unchanged.
At radius zero, evaluation at $U=0$ gives $p(T)$, hence the entire point $x$.
The polynomial $U+T$ has precisely its constant and degree-one coefficients,
so its actual Gauss value is $\max(|T(x)|,s)$. Its constant coefficient after
translation is $p(T)$; the genuine Gauss coefficient bound gives the lower
bound for EVERY original polynomial.

We prove translation invariance on a Gauss disc without imposing any
coefficient binomial bound. If $z$ is any whole coefficient point and
$|a(z)|\leq s$, the translated Gauss point over $z$ has coordinate value
at most $\max(s,|a(z)|)\leq s$, and the same coefficient restriction.
The already proved maximality of the original Gauss point therefore bounds
its value on every polynomial above the translated value. Apply the same
argument to the inverse translation $U\mapsto U-a$ and then to the polynomial
obtained by forward translation. The inverse cancels as an actual algebra
isomorphism, giving the reverse bound on EVERY polynomial. Thus forward
translation fixes the whole Gauss point. Applying this with coefficient
point $x$ and $a=T$ proves that, if $|T(x)|\leq s$, inflation equals the
restriction of the untranslated Gauss point over $x$. Its entire finite
coefficient supremum is precisely the Gauss point of radius $s$ over the
original $A$ restriction. Finally, for $x=\eta_t(y)$ and $s\leq t$, inflation
has coordinate value $t$ and coefficient restriction $y$. Gauss maximality
bounds it above by $x$; the already proved lower bound gives the reverse
inequality on EVERY original polynomial. Whole-point extensionality finishes.
\end{proof}

\begin{theorem}[A genuine strong deformation to the fibre annulus skeleton]
\label{analytic:actual-fibre-annulus-gauss-deformation}

\uses{analytic:actual-polynomial-gauss-inflation}
With the same arbitrary original rings and valuation, fix an actual point
$y\in\operatorname{Sp}_v(A)$ and nonnegative endpoints $\rho,r$. Let
\[
 \mathcal A_y[\rho,r]=\{x\in\operatorname{Sp}_v(A[T]):
 x|_A=y,\ \rho\leq |T(x)|\leq r\}.
\]
The actual radius map $q(x)=|T(x)|$ to the ORIGINAL interval $[\rho,r]$
and the actual Gauss section $j(t)=\eta_t(y)$ are continuous and $qj$ is
exactly the identity. There is a constructed continuous homotopy from the
identity of $\mathcal A_y[\rho,r]$ to $jq$, namely
\[
 H(\theta,x)=I_{\theta |T(x)|}(x),\qquad 0\leq\theta\leq1.
\]
This homotopy preserves the entire coefficient point $y$ and the exact
original coordinate radius. It fixes EVERY point of the Gauss section for
EVERY time, and hence is a genuine strong deformation to its actual image.
No positivity or ordering of the endpoints is needed for the construction;
if $\rho>r$ both the interval and the annulus are actually empty.
\end{theorem}
\begin{proof}
Polynomial evaluation is continuous on the original spectrum, so $q$ is
continuous into the indicated interval subtype. Joint Gauss continuity gives
continuity of $j$, and the actual coordinate value of a Gauss point gives
$q(j(t))=t$. For every allowed $\theta$, one has
$\theta|T(x)|\leq |T(x)|$. The whole coefficient identity and the exact
maximum formula from the previous theorem therefore put $H(\theta,x)$ in
the SAME original fibre annulus and give $q(H(\theta,x))=q(x)$.
Joint continuity of inflation, multiplication, and original coordinate
evaluation proves continuity on $[0,1]\times\mathcal A_y[\rho,r]$ with its
actual subspace topology. At time zero the whole inflation identity gives
$x$. At time one the original coordinate bound gives the entire Gauss point
$\eta_{|T(x)|}(y)=j(q(x))$. For a point $j(t)$, the inflation radius
$\theta t\leq t$ and the whole Gauss fixed-point assertion give
$H(\theta,j(t))=j(t)$. These are actual equalities of all original polynomial
values, not assertions merely about supports or scheme points.
\end{proof}
